% Options for packages loaded elsewhere
\PassOptionsToPackage{unicode}{hyperref}
\PassOptionsToPackage{hyphens}{url}
%
\documentclass[
]{article}
\usepackage{amsmath,amssymb}
\usepackage{iftex}
\ifPDFTeX
  \usepackage[T1]{fontenc}
  \usepackage[utf8]{inputenc}
  \usepackage{textcomp} % provide euro and other symbols
\else % if luatex or xetex
  \usepackage{unicode-math} % this also loads fontspec
  \defaultfontfeatures{Scale=MatchLowercase}
  \defaultfontfeatures[\rmfamily]{Ligatures=TeX,Scale=1}
\fi
\usepackage{lmodern}
\usepackage{float}
\usepackage{tabularx}
\ifPDFTeX\else
  % xetex/luatex font selection
\fi
% Use upquote if available, for straight quotes in verbatim environments
\IfFileExists{upquote.sty}{\usepackage{upquote}}{}
\IfFileExists{microtype.sty}{% use microtype if available
  \usepackage[]{microtype}
  \UseMicrotypeSet[protrusion]{basicmath} % disable protrusion for tt fonts
}{}
\makeatletter
\@ifundefined{KOMAClassName}{% if non-KOMA class
  \IfFileExists{parskip.sty}{%
    \usepackage{parskip}
  }{% else
    \setlength{\parindent}{0pt}
    \setlength{\parskip}{6pt plus 2pt minus 1pt}}
}{% if KOMA class
  \KOMAoptions{parskip=half}}
\makeatother
\usepackage{xcolor}
\usepackage{longtable,booktabs,array}
\usepackage{multirow}
\usepackage{calc} % for calculating minipage widths
% Correct order of tables after \paragraph or \subparagraph
\usepackage{etoolbox}
\makeatletter
\patchcmd\longtable{\par}{\if@noskipsec\mbox{}\fi\par}{}{}
\makeatother
% Allow footnotes in longtable head/foot
\IfFileExists{footnotehyper.sty}{\usepackage{footnotehyper}}{\usepackage{footnote}}
\makesavenoteenv{longtable}
\usepackage{graphicx}
\makeatletter
\def\maxwidth{\ifdim\Gin@nat@width>\linewidth\linewidth\else\Gin@nat@width\fi}
\def\maxheight{\ifdim\Gin@nat@height>\textheight\textheight\else\Gin@nat@height\fi}
\makeatother
% Scale images if necessary, so that they will not overflow the page
% margins by default, and it is still possible to overwrite the defaults
% using explicit options in 

\setkeys{Gin}{width=\maxwidth,height=\maxheight,keepaspectratio}
% Set default figure placement to htbp
\makeatletter
\def\fps@figure{htbp}
\makeatother
\ifLuaTeX
  \usepackage{luacolor}
  \usepackage[soul]{lua-ul}
\else
  \usepackage{soul}
\fi
\setlength{\emergencystretch}{3em} % prevent overfull lines
\providecommand{\tightlist}{%
  \setlength{\itemsep}{0pt}\setlength{\parskip}{0pt}}
\setcounter{secnumdepth}{-\maxdimen} % remove section numbering
\ifLuaTeX
  \usepackage{selnolig}  % disable illegal ligatures
\fi
\IfFileExists{bookmark.sty}{\usepackage{bookmark}}{\usepackage{hyperref}}
\IfFileExists{xurl.sty}{\usepackage{xurl}}{} % add URL line breaks if available
\urlstyle{same}
\hypersetup{
  hidelinks,
  pdfcreator={LaTeX via pandoc}}

\author{}
\date{}

\begin{document}

\begin{longtable}[]{@{}
  >{\raggedright\arraybackslash}p{(\columnwidth - 0\tabcolsep) * \real{1.0000}}@{}}
\toprule\noalign{}
\begin{minipage}[b]{\linewidth}\raggedright

\begin{tabular}{@{}
  >{\raggedright\arraybackslash}p{(\columnwidth - 4\tabcolsep) * \real{0.3333}}
  >{\raggedright\arraybackslash}p{(\columnwidth - 4\tabcolsep) * \real{0.3333}}
  >{\raggedright\arraybackslash}p{(\columnwidth - 4\tabcolsep) * \real{0.3333}}@{}}
\toprule\noalign{}
\multirow{2}{*}{\begin{minipage}[b]{\linewidth}\raggedright

\begin{center}
\fbox{\begin{minipage}{0.9\linewidth}\centering
\textbf{[PLACEHOLDER: Add image here]}
 \textit{Original image path: media/image1.png}
\end{minipage}}
\end{center}

\end{minipage}} & \begin{minipage}[b]{\linewidth}\raggedright
ECOLE SUPERIEUR DES INDUSTRIESDU
\end{minipage} &
\multirow{2}{*}{\begin{minipage}[b]{\linewidth}\raggedright

\begin{center}
\fbox{\begin{minipage}{0.9\linewidth}\centering
\textbf{[PLACEHOLDER: Add image here]}
 \textit{Original image path: media/image2.png}
\end{minipage}}
\end{center}

\end{minipage}} \\
& \begin{minipage}[b]{\linewidth}\raggedright
TEXTILE ET DE L'HABILLEMENT
\end{minipage} & \\
\midrule\noalign{}
\bottomrule\noalign{}
\end{tabular}

\vspace{1em}

\textbf{Rapport de suivit}

\vspace{0.5em}

\begin{tabular}{@{}
  >{\raggedright\arraybackslash}p{(\columnwidth - 2\tabcolsep) * \real{0.5000}}
  >{\raggedright\arraybackslash}p{(\columnwidth - 2\tabcolsep) * \real{0.5000}}@{}}
\toprule\noalign{}
\multicolumn{2}{@{}l@{}}{%
\begin{minipage}[b]{\linewidth}\raggedright
\textbf{Ecole supérieure des industrie textile et habillement}
\end{minipage}} \\
\midrule\noalign{}
Année scolaire & 2025 - 2026 \\
Nom et prénom de l'étudiant : & Mehdi Boumazzourh \\
\bottomrule\noalign{}
\end{tabular}

\vspace{1em}

\textbf{Encadrant de l'ESITH} : Mr SAMIR TETOUANI

\textbf{Domaine} : Lean manufacturing et excellence opérationnelle

\vspace{1em}

\begin{center}
\fbox{\begin{minipage}{0.9\linewidth}\centering
\textbf{[PLACEHOLDER: Add image here]}
 \textit{Original image path: media/image3.png}
\end{minipage}}
\end{center}

\vspace{1em}

Cycle : Génie industriel

Filière : Chef de produit

Niveau : 3\textsuperscript{ème} année

\end{minipage} \\
\midrule\noalign{}
\endhead
\bottomrule\noalign{}
\endlastfoot
\end{longtable}

\begin{longtable}[]{@{}
  >{\raggedright\arraybackslash}p{(\columnwidth - 0\tabcolsep) * \real{1.0000}}@{}}
\toprule\noalign{}
\begin{minipage}[b]{\linewidth}\raggedright

\begin{center}
\fbox{\begin{minipage}{0.9\linewidth}\centering
\textbf{[PLACEHOLDER: Add image here]}
 \textit{Original image path: media/image4.png}
\end{minipage}}
\end{center}

\vspace{0.5em}

\begin{center}
\fbox{\begin{minipage}{0.9\linewidth}\centering
\textbf{[PLACEHOLDER: Add image here]}
 \textit{Original image path: media/image5.png}
\end{minipage}}
\end{center}

\vspace{1em}

\textbf{Définition des termes générales}

\vspace{0.5em}

Le \textbf{Lean} (qui signifie <<~maigre~>> ou <<~sans superflu~>>) est une philosophie de gestion inspirée du système de production de \textbf{Toyota} (TPS). Son objectif principal est de maximiser la valeur pour le client tout en minimisant les gaspillages.

Bien que les deux termes soient souvent utilisés de manière interchangeable, ils désignent des périmètres légèrement différents :

\textbf{1. Le Lean Manufacturing (La dimension opérationnelle)}

C\textquotesingle est l\textquotesingle application historique du Lean au sein des \textbf{usines et des lignes de production}. Il se concentre sur l\textquotesingle optimisation des flux physiques.

\begin{itemize}
\item \textbf{Objectif :} Produire <<~au plus juste~>> (juste ce qu\textquotesingle il faut, quand il faut).
\item \textbf{Action :} Éliminer les \textbf{Muda} (gaspillages) dans l\textquotesingle atelier, comme les stocks excessifs, les déplacements inutiles ou les défauts de fabrication.
\item \textbf{Outils clés :} 5S (organisation), Kanban (flux tirés), SMED (réduction des temps de changement de série), TPM (maintenance). 
\end{itemize}
\textbf{2. Le Lean Management (La dimension organisationnelle)}

Il s\textquotesingle agit de l\textquotesingle évolution du Lean Manufacturing à \textbf{l\textquotesingle ensemble de l\textquotesingle entreprise} (bureaux, RH, logistique, service client, direction). C\textquotesingle est autant un système de gestion qu\textquotesingle une culture d\textquotesingle entreprise.

\begin{itemize}
\item \textbf{Objectif :} Créer une culture d\textquotesingle amélioration continue (\textbf{Kaizen}) impliquant tous les collaborateurs.
\item \textbf{Action :} Optimiser les processus administratifs et décisionnels, et surtout, développer les compétences des employés pour qu\textquotesingle ils résolvent eux-mêmes les problèmes.
\item \textbf{Focus :} Le leadership, l\textquotesingle écoute du client et l\textquotesingle intelligence collective. 
\end{itemize}
\vspace{1em}

\textbf{Tableau Comparatif Rapide}

\vspace{0.5em}

\begin{tabular}{@{}
  >{\raggedright\arraybackslash}p{0.25\linewidth}
  >{\raggedright\arraybackslash}p{0.35\linewidth}
  >{\raggedright\arraybackslash}p{0.35\linewidth}@{}}
\toprule
\textbf{Caractéristique} & \textbf{Lean Manufacturing} & \textbf{Lean Management} \\
\midrule
\textbf{Périmètre} & L\textquotesingle atelier / La production & L\textquotesingle organisation entière \\
\midrule
\textbf{Cible} & Les machines, les flux, les pièces & Les personnes, les processus, la stratégie \\
\midrule
\textbf{Outils types} & VSM, 5S, Kanban, Jidoka & Management Visuel, Kaizen, A3, Hoshin Kanri \\
\midrule
\textbf{Résultats visés} & Gain de productivité, réduction de stock & Agilité, satisfaction client, engagement salarié \\
\bottomrule
\end{tabular}

\end{minipage} \\
\midrule\noalign{}
\endhead
\bottomrule\noalign{}
\endlastfoot
\end{longtable}

\begin{longtable}[]{@{}
  >{\raggedright\arraybackslash}p{(\columnwidth - 0\tabcolsep) * \real{1.0000}}@{}}
\toprule\noalign{}
\begin{minipage}[b]{\linewidth}\raggedright
\begin{center}
\fbox{\begin{minipage}{0.9\linewidth}\centering
\textbf{[PLACEHOLDER: Add image here]}
 \textit{Original image path: media/image4.png}
\end{minipage}}
\end{center}

\begin{center}
\fbox{\begin{minipage}{0.9\linewidth}\centering
\textbf{[PLACEHOLDER: Add image here]}
 \textit{Original image path: media/image5.png}
\end{minipage}}
\end{center}

\textbf{Les 5 principes fondamentaux du Lean}

Pour réussir dans l\textquotesingle une ou l\textquotesingle autre de
ces approches, on suit généralement ces 5 étapes : 1.\textbf{Définir la
Valeur :} Qu\textquotesingle est-ce que le client est réellement prêt à
payer ?

2.\textbf{Identifier le Flux de Valeur (VSM) :} Cartographier toutes les
étapes pour repérer celles qui ne servent à rien.

3.\textbf{Créer un Flux Continu :} Faire en sorte que le produit/service
avance sans interruption (ni attente, ni stock).

4.\textbf{Mettre en place le Flux Tiré :} Ne produire que
lorsqu\textquotesingle il y a une demande réelle du client.
5.\textbf{Viser la Perfection :} Répéter le cycle pour
s\textquotesingle améliorer sans cesse (Kaizen).

Comparaison entre excellence opération et lean management et origine des
deux

\begin{tabular}{@{}
  >{\raggedright\arraybackslash}p{0.25\linewidth}
  >{\raggedright\arraybackslash}p{0.35\linewidth}
  >{\raggedright\arraybackslash}p{0.35\linewidth}@{}}
\toprule
\textbf{Caractéristique} & \textbf{Lean Management} & \textbf{Excellence Opérationnelle} \\
\midrule
\textbf{Définition} & Philosophie visant à éliminer les gaspillages pour maximiser la valeur client. & Démarche systémique visant une performance durable et une culture d\textquotesingle amélioration. \\
\midrule
\textbf{Focus principal} & Flux, vitesse et réduction des gaspillages (Muda). & Culture, leadership, stratégie et résultats globaux. \\
\midrule
\textbf{Boîte à outils} & 5S, VSM, Kanban, Kaizen, SMED. & Lean + Six Sigma + TPM + Management stratégique + Digital. \\
\midrule
\textbf{Vision} & <<~Comment faire mieux avec moins~?~>> & <<~Comment ancrer la performance dans l\textquotesingle ADN de l\textquotesingle entreprise~?~>> \\
\midrule
\textbf{Acteurs} & Souvent perçu comme une démarche terrain (Gemba). & Implication totale, de la direction générale aux opérateurs. \\
\bottomrule
\end{tabular}

\vspace{1em}

\textbf{Définition des termes et de quelques méthodes}

\textbf{1. Définition des Termes et Méthodes Clés}

Avant d\textquotesingle aborder les problèmes, clarifions le vocabulaire technique :

\begin{itemize}
\item \textbf{Lean Manufacturing :} Philosophie visant à éliminer les \textbf{Muda} (gaspillages) pour maximiser la valeur ajoutée avec le moins de ressources possible.
\item \textbf{Six Sigma :} Méthodologie statistique visant à réduire la \textbf{variabilité} des processus pour atteindre un taux de défaut quasi nul (3,4 défauts par million d\textquotesingle opportunités).
\item \textbf{Lean Six Sigma (LSS) :} L\textquotesingle hybride des deux : la vitesse et l\textquotesingle agilité du Lean combinées à la rigueur statistique du Six Sigma. 
\end{itemize}
\end{minipage} \\
\midrule\noalign{}
\endhead
\bottomrule\noalign{}
\endlastfoot
\end{longtable}

\begin{longtable}[]{@{}
  >{\raggedright\arraybackslash}p{(\columnwidth - 0\tabcolsep) * \real{1.0000}}@{}}
\toprule\noalign{}
\begin{minipage}[b]{\linewidth}\raggedright
\begin{center}
\fbox{\begin{minipage}{0.9\linewidth}\centering
\textbf{[PLACEHOLDER: Add image here]}
 \textit{Original image path: media/image4.png}
\end{minipage}}
\end{center}

\begin{center}
\fbox{\begin{minipage}{0.9\linewidth}\centering
\textbf{[PLACEHOLDER: Add image here]}
 \textit{Original image path: media/image5.png}
\end{minipage}}
\end{center}

\begin{itemize}
\item \textbf{Just-In-Time (JIT) :} Produire exactement ce qui est nécessaire, quand c\textquotesingle est nécessaire et en quantité
\end{itemize}
voulue. Repose sur le flux tiré (\textbf{Kanban}).

\begin{itemize}
\item \textbf{Kaizen :} Démarche d\textquotesingle amélioration continue par "petits pas", impliquant tout le personnel, du terrain à la direction.
\end{itemize}

\begin{itemize}
\item \textbf{VSM (Value Stream Mapping) :} Cartographie du flux de valeur permettant de visualiser les gaspillages entre le fournisseur et le
\end{itemize}
client.

\textbf{2. Typologie des Problèmes Industriels}

En industrie, on classe généralement les problèmes selon
l\textquotesingle indicateur qu\textquotesingle ils dégradent (modèle
\textbf{SQCDP}) :

1.\textbf{Qualité (Quality) :} Taux de rebuts trop élevé,
non-conformités, réclamations clients, variabilité dimensionnelle.

2.\textbf{Productivité/Coût (Cost) :} Temps de cycle trop long, stocks
excessifs (WIP), pannes machines fréquentes, sous-optimisation de la
main-d\textquotesingle œuvre.

3.\textbf{Délai (Delivery) :} Retards de livraison, goulots
d\textquotesingle étranglement dans le flux, manque de composants
(ruptures).

4.\textbf{Sécurité \& Environnement (Safety/Environment) :} Risques
d\textquotesingle accidents, ergonomie des postes (TMS), consommation
énergétique excessive.

\begin{center}
\fbox{\begin{minipage}{0.9\linewidth}\centering
\textbf{[PLACEHOLDER: Add image here]}
 \textit{Original image path: media/image6.png}
\end{minipage}}
\end{center}

DMAIC
\end{minipage} \\
\midrule\noalign{}
\endhead
\bottomrule\noalign{}
\endlastfoot
\end{longtable}

\begin{longtable}[]{@{}
  >{\raggedright\arraybackslash}p{(\columnwidth - 0\tabcolsep) * \real{1.0000}}@{}}
\toprule\noalign{}
\begin{minipage}[b]{\linewidth}\raggedright
\begin{center}
\fbox{\begin{minipage}{0.9\linewidth}\centering
\textbf{[PLACEHOLDER: Add image here]}
 \textit{Original image path: media/image4.png}
\end{minipage}}
\end{center}

\begin{center}
\fbox{\begin{minipage}{0.9\linewidth}\centering
\textbf{[PLACEHOLDER: Add image here]}
 \textit{Original image path: media/image5.png}
\end{minipage}}
\end{center}

En milieu industriel, on distingue ces deux situations par la nature de
l\textquotesingle{}\textbf{écart} (le "Gap") :

\textbf{1. La 8D : La gestion de l\textquotesingle Anomalie
(Restauration)}

\begin{itemize}
\item \textbf{Situation :} Le standard existe, il était maîtrisé, mais "quelque chose" a cassé le processus.
\end{itemize}

\begin{itemize}
\item \textbf{L\textquotesingle objectif :} C\textquotesingle est un retour vers le passé (revenir à l\textquotesingle état stable que
\end{itemize}
l\textquotesingle on connaissait).

\begin{itemize}
\item \textbf{Le mot-clé :Réactivité}. On parle de "Correctif".

\item \textbf{Analogie :} Un athlète qui court habituellement le 100m en 10 secondes (objectif atteint) et
\end{itemize}

qui, suite à une blessure, tombe à 15 secondes. La 8D,
c\textquotesingle est le protocole de soin pour qu\textquotesingle il

court à nouveau en 10 secondes.

\textbf{2. Le DMAIC : La gestion du Projet (Rupture)}

\begin{itemize}
\item \textbf{Situation :} Le processus est stable, il n\textquotesingle y a pas de "panne", mais le niveau actuel ne suffit plus
\end{itemize}

pour être compétitif.

\begin{itemize}
\item \textbf{L\textquotesingle objectif :} C\textquotesingle est un bond vers le futur (atteindre un niveau de performance jamais atteint
\end{itemize}

auparavant).

\begin{itemize}
\item \textbf{Le mot-clé :Proactif}. On parle d\textquotesingle{} Amélioration de rupture (Breakthrough).
\end{itemize}

\begin{itemize}
\item \textbf{Analogie :} Le même athlète court toujours en 10 secondes (stable), mais pour gagner les JO,
\end{itemize}

il doit passer à 9,7 secondes. Il ne s\textquotesingle agit pas de
soigner une blessure, mais de changer de

régime, de technique de course et d\textquotesingle équipement.
C\textquotesingle est le DMAIC.

Situation,Question à se poser,Méthode

Chute soudaine → Qu'est-ce qui a changé par rapport à hier ? → ,8D

Objectif futur, →Comment faire mieux que ce que l\textquotesingle on
sait déjà faire ? →, DMAIC
\end{minipage} \\
\midrule\noalign{}
\endhead
\bottomrule\noalign{}
\endlastfoot
\end{longtable}

\begin{longtable}[]{@{}
  >{\raggedright\arraybackslash}p{(\columnwidth - 0\tabcolsep) * \real{1.0000}}@{}}
\toprule\noalign{}
\begin{minipage}[b]{\linewidth}\raggedright
\begin{center}
\fbox{\begin{minipage}{0.9\linewidth}\centering
\textbf{[PLACEHOLDER: Add image here]}
 \textit{Original image path: media/image4.png}
\end{minipage}}
\end{center}

\begin{center}
\fbox{\begin{minipage}{0.9\linewidth}\centering
\textbf{[PLACEHOLDER: Add image here]}
 \textit{Original image path: media/image5.png}
\end{minipage}}
\end{center}

\textbf{Etude de cas}

\textbf{Définir}

Problématique :

\begin{tabular}{@{}
  >{\raggedright\arraybackslash}p{(\columnwidth - 2\tabcolsep) * \real{0.5000}}
  >{\raggedright\arraybackslash}p{(\columnwidth - 2\tabcolsep) * \real{0.5000}}@{}}
\toprule\noalign{}
\multicolumn{2}{@{}>{\raggedright\arraybackslash}p{(\columnwidth - 2\tabcolsep) * \real{1.0000} + 2\tabcolsep}@{}}{%
\begin{minipage}[b]{\linewidth}\raggedright
Section Détails du Projet (Données Robotica)
\end{minipage}} \\
\midrule\noalign{}

\bottomrule\noalign{}

\begin{minipage}[t]{\linewidth}\raggedright
Situation actuelle
\end{minipage} & \begin{minipage}[t]{\linewidth}\raggedright
\begin{itemize}
\item Mécontentement des donneurs d'ordres : retards de livraison et coût élevé des défauts qualité.
\item Menace de changement de fournisseur sous 6 mois si la situation n'est pas redressée.
\item Performance technique : 45 pannes/mois, temps de changement de référence de 3h, délai de livraison moyen de 7 jours.
\end{itemize}
\end{minipage} \\
\begin{minipage}[t]{\linewidth}\raggedright
Situation future
\end{minipage} & \begin{minipage}[t]{\linewidth}\raggedright
\begin{itemize}
\item Clients satisfaits et fidélisés grâce à une meilleure fiabilité.
\item Délais de livraison réduits de 20\% (cible : 5,6 jours).
\item Entreprise rentable avec une marge supérieure à 20\%.
\item Processus stabilisé avec une réduction significative des pannes et du temps SMED.
\end{itemize}
\end{minipage} \\
\begin{minipage}[t]{\linewidth}\raggedright
Objectif
\end{minipage} & \begin{minipage}[t]{\linewidth}\raggedright
\begin{itemize}
\item Réduction du coût de revient de 15\%.
\item Réduction des délais de livraison de 20\%.
\item Augmentation de la marge à plus de 20\%.
\end{itemize}
\end{minipage} \\
\begin{minipage}[t]{\linewidth}\raggedright
Équipe de projet
\end{minipage} & \begin{minipage}[t]{\linewidth}\raggedright
\begin{itemize}
\item Sponsor : Directeur Général (DG) de Robotica.

\item Leader : Responsable Excellence Opérationnelle (EO).

\item Partenaires : Équipes R\&D (Allemagne) et Production/Montage (France). \end{itemize}
\end{minipage} \\
\begin{minipage}[t]{\linewidth}\raggedright
macro planning
\end{minipage} & \begin{minipage}[t]{\linewidth}\raggedright
\begin{itemize}
\item Remise du rapport de diagnostic : sous 20 heures.

\item Phase de redressement globale : maximum 6 mois. \end{itemize}
\end{minipage} \\
\begin{minipage}[t]{\linewidth}\raggedright
risque de projet
\end{minipage} & \begin{minipage}[t]{\linewidth}\raggedright
\begin{itemize}
\item Perte définitive des clients majeurs si le délai de 6 mois n'est pas respecté.
\end{itemize}

\begin{itemize}
\item Instabilité critique de la production due au nombre élevé de pannes (45/mois).
\end{itemize}

\begin{itemize}
\item Inflexibilité face à la variation de la demande (60 à 100 commandes). \end{itemize}
\end{minipage} \\
Plan de commmunication & \begin{minipage}[t]{\linewidth}\raggedright
\begin{itemize}
\item Reporting direct et immédiat au DG (diagnostic sous 20h).

\item Coordination inter-sites entre la R\&D allemande et l'usine française.

\item Suivi mensuel des indicateurs clés (KPI) de performance pour les donneurs
\end{itemize}

d'ordres.
\end{minipage} \\
\end{tabular}

\hl{Note du cours :}

\hl{Charte de projet DMAIC :}

\begin{tabular}{@{}
  >{\raggedright\arraybackslash}p{(\columnwidth - 4\tabcolsep) * \real{0.3333}}
  >{\raggedright\arraybackslash}p{(\columnwidth - 4\tabcolsep) * \real{0.3333}}
  >{\raggedright\arraybackslash}p{(\columnwidth - 4\tabcolsep) * \real{0.3333}}@{}}
\toprule\noalign{}
\begin{minipage}[b]{\linewidth}\raggedright
~
\end{minipage} &
\multicolumn{2}{>{\raggedright\arraybackslash}p{(\columnwidth - 4\tabcolsep) * \real{0.6667} + 2\tabcolsep}@{}}{%
\begin{minipage}[b]{\linewidth}\raggedright
Equipe de projet :
\end{minipage}} \\
\midrule\noalign{}

\bottomrule\noalign{}

& 
\multicolumn{2}{>{\raggedright\arraybackslash}p{(\columnwidth - 4\tabcolsep) * \real{0.6667} + 2\tabcolsep}@{}}{%
\begin{minipage}[t]{\linewidth}\raggedright
Macro-planning
\end{minipage}} \\
& &  \begin{minipage}[t]{\linewidth}\raggedright
D \textbar-\/-2S-\/-\textbar{}
\end{minipage} \\
& &  \begin{minipage}[t]{\linewidth}\raggedright
M \textbar-\/-3M-\/-\textbar{}
\end{minipage} \\
& &  \begin{minipage}[t]{\linewidth}\raggedright
A\ldots\ldots\ldots\ldots\ldots{}
\end{minipage} \\
& &  \begin{minipage}[t]{\linewidth}\raggedright
I\ldots\ldots\ldots\ldots{}
\end{minipage} \\
& &  \begin{minipage}[t]{\linewidth}\raggedright
C\ldots\ldots\ldots.
\end{minipage} \\
& 
\multicolumn{2}{>{\raggedright\arraybackslash}p{(\columnwidth - 4\tabcolsep) * \real{0.6667} + 2\tabcolsep}@{}}{%
\begin{minipage}[t]{\linewidth}\raggedright
Risque du projet :
\end{minipage}} \\
\end{tabular}
\end{minipage} \\
\midrule\noalign{}
\endhead
\bottomrule\noalign{}
\endlastfoot
\end{longtable}

\begin{center}
\fbox{\begin{minipage}{0.9\linewidth}\centering
\textbf{[PLACEHOLDER: Add image here]}
 \textit{Original image path: media/image7.png}
\end{minipage}}
\end{center}

\begin{center}
\fbox{\begin{minipage}{0.9\linewidth}\centering
\textbf{[PLACEHOLDER: Add image here]}
 \textit{Original image path: media/image8.png}
\end{minipage}}
\end{center}

\begin{center}
\fbox{\begin{minipage}{0.9\linewidth}\centering
\textbf{[PLACEHOLDER: Add image here]}
 \textit{Original image path: media/image9.png}
\end{minipage}}
\end{center}

\begin{center}
\fbox{\begin{minipage}{0.9\linewidth}\centering
\textbf{[PLACEHOLDER: Add image here]}
 \textit{Original image path: media/image10.png}
\end{minipage}}
\end{center}

\begin{center}
\fbox{\begin{minipage}{0.9\linewidth}\centering
\textbf{[PLACEHOLDER: Add image here]}
 \textit{Original image path: media/image11.png}
\end{minipage}}
\end{center}

\begin{center}
\fbox{\begin{minipage}{0.9\linewidth}\centering
\textbf{[PLACEHOLDER: Add image here]}
 \textit{Original image path: media/image12.png}
\end{minipage}}
\end{center}

\begin{center}
\fbox{\begin{minipage}{0.9\linewidth}\centering
\textbf{[PLACEHOLDER: Add image here]}
 \textit{Original image path: media/image13.png}
\end{minipage}}
\end{center}

\begin{center}
\fbox{\begin{minipage}{0.9\linewidth}\centering
\textbf{[PLACEHOLDER: Add image here]}
 \textit{Original image path: media/image4.png}
\end{minipage}}
\end{center}

\begin{center}
\fbox{\begin{minipage}{0.9\linewidth}\centering
\textbf{[PLACEHOLDER: Add image here]}
 \textit{Original image path: media/image14.png}
\end{minipage}}
\end{center}

\begin{center}
\fbox{\begin{minipage}{0.9\linewidth}\centering
\textbf{[PLACEHOLDER: Add image here]}
 \textit{Original image path: media/image15.png}
\end{minipage}}
\end{center}

\begin{center}
\fbox{\begin{minipage}{0.9\linewidth}\centering
\textbf{[PLACEHOLDER: Add image here]}
 \textit{Original image path: media/image16.png}
\end{minipage}}
\end{center}

\begin{center}
\fbox{\begin{minipage}{0.9\linewidth}\centering
\textbf{[PLACEHOLDER: Add image here]}
 \textit{Original image path: media/image17.png}
\end{minipage}}
\end{center}

\begin{center}
\fbox{\begin{minipage}{0.9\linewidth}\centering
\textbf{[PLACEHOLDER: Add image here]}
 \textit{Original image path: media/image18.png}
\end{minipage}}
\end{center}

\begin{center}
\fbox{\begin{minipage}{0.9\linewidth}\centering
\textbf{[PLACEHOLDER: Add image here]}
 \textit{Original image path: media/image19.png}
\end{minipage}}
\end{center}

\begin{center}
\fbox{\begin{minipage}{0.9\linewidth}\centering
\textbf{[PLACEHOLDER: Add image here]}
 \textit{Original image path: media/image5.png}
\end{minipage}}
\end{center}

\begin{center}
\fbox{\begin{minipage}{0.9\linewidth}\centering
\textbf{[PLACEHOLDER: Add image here]}
 \textit{Original image path: media/image20.png}
\end{minipage}}
\end{center}

\begin{center}
\fbox{\begin{minipage}{0.9\linewidth}\centering
\textbf{[PLACEHOLDER: Add image here]}
 \textit{Original image path: media/image21.png}
\end{minipage}}
\end{center}

\begin{center}
\fbox{\begin{minipage}{0.9\linewidth}\centering
\textbf{[PLACEHOLDER: Add image here]}
 \textit{Original image path: media/image22.png}
\end{minipage}}
\end{center}

\begin{longtable}[]{@{}
  >{\raggedright\arraybackslash}p{(\columnwidth - 0\tabcolsep) * \real{1.0000}}@{}}
\toprule\noalign{}
\begin{minipage}[b]{\linewidth}\raggedright
\begin{itemize}
\item Plan de communication :\\ \textbf{Mesurer} Taux de client
\end{itemize}

Mesure de Kpi's :\\
\hl{Taux client → objectif}

Objectif \hl{Voix du client :} CTQ( critical to quality ) =KPI\\
Voix du proess : Taux de satisfaction client :\\
\begin{itemize}
\item KPI délai\\
\item \\ Cout • Taux de conformité
\end{itemize}

\textbf{Synthèse Globale}

Le total du Chiffre d\textquotesingle Affaires (CA) est de 60 386 360
pour une Quantité Totale vendue de 528 161 unités, réparties sur 1 456
transactions. La base de données couvre 10 produits uniques et 5 clients
uniques.

\textbf{1. Prise de Vue Générale et Concentration des Revenus}

L\textquotesingle analyse révèle une forte concentration du Chiffre
d\textquotesingle Affaires, tant au niveau des produits
qu\textquotesingle au niveau des clients.

\begin{itemize}
\item \textbf{Concentration des Produits :} Les \textbf{5 premiers produits} (sur 10) génèrent \textbf{67.66\%} du Chiffre d\textquotesingle Affaires
\end{itemize}
total.

o Les produits \textbf{P02} et \textbf{P10} sont les principaux moteurs
de revenus, suivis par \textbf{P04}, \textbf{P07}, et \textbf{P08}.

o Les produits \textbf{P02} et \textbf{P10} sont également les produits
les plus vendus en volume (quantité), avec des volumes très élevés par
rapport aux autres produits du Top 5, ce qui suggère
qu\textquotesingle ils sont des articles à forte rotation ou à volume
élevé.

\begin{itemize}
\item \textbf{Concentration des Clients :} Étant donné qu\textquotesingle il n\textquotesingle y a que 5 clients uniques, le top 5 des clients génère
\end{itemize}
logiquement \textbf{100.00\%} du CA total.

o Les trois principaux clients sont \textbf{Staübli}, \textbf{Kawasaki}
et \textbf{ABB}, qui contribuent ensemble à la majeure partie du Chiffre
d\textquotesingle Affaires.
\end{minipage} \\
\midrule\noalign{}
\endhead
\bottomrule\noalign{}
\endlastfoot
\end{longtable}

\begin{longtable}[]{@{}
  >{\raggedright\arraybackslash}p{(\columnwidth - 0\tabcolsep) * \real{1.0000}}@{}}
\toprule\noalign{}
\begin{minipage}[b]{\linewidth}\raggedright
\begin{center}
\fbox{\begin{minipage}{0.9\linewidth}\centering
\textbf{[PLACEHOLDER: Add image here]}
 \textit{Original image path: media/image4.png}
\end{minipage}}
\end{center}

\begin{center}
\fbox{\begin{minipage}{0.9\linewidth}\centering
\textbf{[PLACEHOLDER: Add image here]}
 \textit{Original image path: media/image5.png}
\end{minipage}}
\end{center}

\textbf{2. Performance Détaillée}

\begin{tabular}{@{}
  >{\raggedright\arraybackslash}p{(\columnwidth - 0\tabcolsep) * \real{1.0000}}@{}}
\toprule\noalign{}
\begin{minipage}[b]{\linewidth}\raggedright
Somme de QTE par CLIENT 450

\begin{tabular}{@{}
  >{\raggedright\arraybackslash}p{(\columnwidth - 6\tabcolsep) * \real{0.2500}}
  >{\raggedright\arraybackslash}p{(\columnwidth - 6\tabcolsep) * \real{0.2500}}
  >{\raggedright\arraybackslash}p{(\columnwidth - 6\tabcolsep) * \real{0.2500}}
  >{\raggedright\arraybackslash}p{(\columnwidth - 6\tabcolsep) * \real{0.2500}}@{}}
\toprule\noalign{}
\begin{minipage}[b]{\linewidth}\raggedright
400
\end{minipage} & \begin{minipage}[b]{\linewidth}\raggedright
360
\end{minipage} & \begin{minipage}[b]{\linewidth}\raggedright
\end{minipage} & \begin{minipage}[b]{\linewidth}\raggedright
388
\end{minipage} \\
\midrule\noalign{}

\bottomrule\noalign{}

\begin{minipage}[t]{\linewidth}\raggedright
350
\end{minipage} & & 354 & \\
\end{tabular}

300\\
250

\begin{tabular}{@{}
  >{\raggedright\arraybackslash}p{(\columnwidth - 4\tabcolsep) * \real{0.3333}}
  >{\raggedright\arraybackslash}p{(\columnwidth - 4\tabcolsep) * \real{0.3333}}
  >{\raggedright\arraybackslash}p{(\columnwidth - 4\tabcolsep) * \real{0.3333}}@{}}
\toprule\noalign{}
\begin{minipage}[b]{\linewidth}\raggedright
200
\end{minipage} & \begin{minipage}[b]{\linewidth}\raggedright
184
\end{minipage} & \begin{minipage}[b]{\linewidth}\raggedright
170
\end{minipage} \\
\midrule\noalign{}

\bottomrule\noalign{}

\end{tabular}

150\\
100\\
50\\
0

\begin{tabular}{@{}
  >{\raggedright\arraybackslash}p{(\columnwidth - 8\tabcolsep) * \real{0.2000}}
  >{\raggedright\arraybackslash}p{(\columnwidth - 8\tabcolsep) * \real{0.2000}}
  >{\raggedright\arraybackslash}p{(\columnwidth - 8\tabcolsep) * \real{0.2000}}
  >{\raggedright\arraybackslash}p{(\columnwidth - 8\tabcolsep) * \real{0.2000}}
  >{\raggedright\arraybackslash}p{(\columnwidth - 8\tabcolsep) * \real{0.2000}}@{}}
\toprule\noalign{}
\begin{minipage}[b]{\linewidth}\raggedright
ABB
\end{minipage} & \begin{minipage}[b]{\linewidth}\raggedright
Fanuc
\end{minipage} & \begin{minipage}[b]{\linewidth}\raggedright
Kawasaki
\end{minipage} & \begin{minipage}[b]{\linewidth}\raggedright
Kuka
\end{minipage} & \begin{minipage}[b]{\linewidth}\raggedright
Staübli
\end{minipage} \\
\midrule\noalign{}

\bottomrule\noalign{}

\end{tabular}
\end{minipage} \\
\midrule\noalign{}

\bottomrule\noalign{}

\end{tabular}

\emph{Figure 1 Somme de QTE par Client}

Le graphique montre une concentration de l\textquotesingle activité sur
trois clients principaux. \textbf{Staübli} est le client le plus
important, tandis que \textbf{Kuka} et \textbf{Fanuc} semblent être,
dans ce contexte précis, des clients secondaires en termes de volume
(représentant environ 25\% du volume total combiné, contre plus de 75\%
pour les trois autres).

\begin{tabular}{@{}
  >{\raggedright\arraybackslash}p{(\columnwidth - 0\tabcolsep) * \real{1.0000}}@{}}
\toprule\noalign{}
\begin{minipage}[b]{\linewidth}\raggedright
Total

\begin{tabular}{@{}
  >{\raggedright\arraybackslash}p{(\columnwidth - 4\tabcolsep) * \real{0.3333}}
  >{\raggedright\arraybackslash}p{(\columnwidth - 4\tabcolsep) * \real{0.3333}}
  >{\raggedright\arraybackslash}p{(\columnwidth - 4\tabcolsep) * \real{0.3333}}@{}}
\toprule\noalign{}
\begin{minipage}[b]{\linewidth}\raggedright
300000
\end{minipage} & \begin{minipage}[b]{\linewidth}\raggedright
261184
\end{minipage} & \begin{minipage}[b]{\linewidth}\raggedright
254320
\end{minipage} \\
\midrule\noalign{}

\bottomrule\noalign{}

\end{tabular}

250000

200000

\begin{tabular}{@{}
  >{\raggedright\arraybackslash}p{(\columnwidth - 2\tabcolsep) * \real{0.5000}}
  >{\raggedright\arraybackslash}p{(\columnwidth - 2\tabcolsep) * \real{0.5000}}@{}}
\toprule\noalign{}
\begin{minipage}[b]{\linewidth}\raggedright
150000
\end{minipage} & \begin{minipage}[b]{\linewidth}\raggedright
Total
\end{minipage} \\
\midrule\noalign{}

\bottomrule\noalign{}

\end{tabular}

100000

50000

\begin{tabular}{@{}
  >{\raggedright\arraybackslash}p{(\columnwidth - 20\tabcolsep) * \real{0.0909}}
  >{\raggedright\arraybackslash}p{(\columnwidth - 20\tabcolsep) * \real{0.0909}}
  >{\raggedright\arraybackslash}p{(\columnwidth - 20\tabcolsep) * \real{0.0909}}
  >{\raggedright\arraybackslash}p{(\columnwidth - 20\tabcolsep) * \real{0.0909}}
  >{\raggedright\arraybackslash}p{(\columnwidth - 20\tabcolsep) * \real{0.0909}}
  >{\raggedright\arraybackslash}p{(\columnwidth - 20\tabcolsep) * \real{0.0909}}
  >{\raggedright\arraybackslash}p{(\columnwidth - 20\tabcolsep) * \real{0.0909}}
  >{\raggedright\arraybackslash}p{(\columnwidth - 20\tabcolsep) * \real{0.0909}}
  >{\raggedright\arraybackslash}p{(\columnwidth - 20\tabcolsep) * \real{0.0909}}
  >{\raggedright\arraybackslash}p{(\columnwidth - 20\tabcolsep) * \real{0.0909}}
  >{\raggedright\arraybackslash}p{(\columnwidth - 20\tabcolsep) * \real{0.0909}}@{}}
\toprule\noalign{}
\begin{minipage}[b]{\linewidth}\raggedright
0
\end{minipage} & \begin{minipage}[b]{\linewidth}\raggedright
107
\end{minipage} & \begin{minipage}[b]{\linewidth}\raggedright
\end{minipage} & \begin{minipage}[b]{\linewidth}\raggedright
143
\end{minipage} & \begin{minipage}[b]{\linewidth}\raggedright
4247
\end{minipage} & \begin{minipage}[b]{\linewidth}\raggedright
145
\end{minipage} & \begin{minipage}[b]{\linewidth}\raggedright
212
\end{minipage} & \begin{minipage}[b]{\linewidth}\raggedright
4010
\end{minipage} & \begin{minipage}[b]{\linewidth}\raggedright
3649
\end{minipage} & \begin{minipage}[b]{\linewidth}\raggedright
144
\end{minipage} & \begin{minipage}[b]{\linewidth}\raggedright
\end{minipage} \\
\midrule\noalign{}

\bottomrule\noalign{}

& P01 & P02 & P03 & P04 & P05 & P06 & P07 & P08 & P09 & P10 \\
\end{tabular} \\
\midrule\noalign{}

\bottomrule\noalign{}

\end{tabular}

D\textquotesingle un point de vue de gestionnaire industriel ou de data
analyst, cette répartition impose des stratégies différenciées :

\begin{itemize}
\item \textbf{Gestion des Stocks :} Pour \textbf{P02} et \textbf{P10}, une rupture de stock serait catastrophique. Ils nécessitent un suivi
\end{itemize}
rigoureux (inventaire permanent, stock de sécurité optimisé). Pour les
autres, une gestion sur commande ou un stock minimal suffit.
\end{minipage} \\
\midrule\noalign{}
\endhead
\bottomrule\noalign{}
\endlastfoot
\end{longtable}

\begin{longtable}[]{@{}
  >{\raggedright\arraybackslash}p{(\columnwidth - 0\tabcolsep) * \real{1.0000}}@{}}
\toprule\noalign{}
\begin{minipage}[b]{\linewidth}\raggedright
\begin{center}
\fbox{\begin{minipage}{0.9\linewidth}\centering
\textbf{[PLACEHOLDER: Add image here]}
 \textit{Original image path: media/image4.png}
\end{minipage}}
\end{center}
\begin{itemize}
\item \textbf{Optimisation de la Production :} Il est probable que vos lignes de production soient
\end{itemize}

\begin{center}
\fbox{\begin{minipage}{0.9\linewidth}\centering
\textbf{[PLACEHOLDER: Add image here]}
 \textit{Original image path: media/image5.png}
\end{minipage}}
\end{center}

quasi exclusivement configurées pour P02 et P10. Le passage à la
production des autres références (comme P01) génère probablement des
"temps de changement de série" (SMED) coûteux par rapport au volume
produit.

\begin{itemize}
\item \textbf{Analyse de Rentabilité :} Il serait intéressant de vérifier si ces deux produits à fort volume sont aussi ceux qui génèrent le plus de
\end{itemize}
marge, ou s\textquotesingle ils servent de "produits
d\textquotesingle appel".

\begin{tabular}{@{}
  >{\raggedright\arraybackslash}p{(\columnwidth - 0\tabcolsep) * \real{1.0000}}@{}}
\toprule\noalign{}
\begin{minipage}[b]{\linewidth}\raggedright
Total

\begin{tabular}{@{}
  >{\raggedright\arraybackslash}p{(\columnwidth - 4\tabcolsep) * \real{0.3333}}
  >{\raggedright\arraybackslash}p{(\columnwidth - 4\tabcolsep) * \real{0.3333}}
  >{\raggedright\arraybackslash}p{(\columnwidth - 4\tabcolsep) * \real{0.3333}}@{}}
\toprule\noalign{}
\begin{minipage}[b]{\linewidth}\raggedright
12000000
\end{minipage} & \begin{minipage}[b]{\linewidth}\raggedright
10447360
\end{minipage} & \begin{minipage}[b]{\linewidth}\raggedright
10172800
\end{minipage} \\
\midrule\noalign{}

\bottomrule\noalign{}

\end{tabular}

10000000

\begin{tabular}{@{}
  >{\raggedright\arraybackslash}p{(\columnwidth - 18\tabcolsep) * \real{0.1000}}
  >{\raggedright\arraybackslash}p{(\columnwidth - 18\tabcolsep) * \real{0.1000}}
  >{\raggedright\arraybackslash}p{(\columnwidth - 18\tabcolsep) * \real{0.1000}}
  >{\raggedright\arraybackslash}p{(\columnwidth - 18\tabcolsep) * \real{0.1000}}
  >{\raggedright\arraybackslash}p{(\columnwidth - 18\tabcolsep) * \real{0.1000}}
  >{\raggedright\arraybackslash}p{(\columnwidth - 18\tabcolsep) * \real{0.1000}}
  >{\raggedright\arraybackslash}p{(\columnwidth - 18\tabcolsep) * \real{0.1000}}
  >{\raggedright\arraybackslash}p{(\columnwidth - 18\tabcolsep) * \real{0.1000}}
  >{\raggedright\arraybackslash}p{(\columnwidth - 18\tabcolsep) * \real{0.1000}}
  >{\raggedright\arraybackslash}p{(\columnwidth - 18\tabcolsep) * \real{0.1000}}@{}}
\toprule\noalign{}
\begin{minipage}[b]{\linewidth}\raggedright
8000000
\end{minipage} & \begin{minipage}[b]{\linewidth}\raggedright
\end{minipage} & \begin{minipage}[b]{\linewidth}\raggedright
\end{minipage} &
\multirow{2}{*}{\begin{minipage}[b]{\linewidth}\raggedright
7219900
\end{minipage}} & \begin{minipage}[b]{\linewidth}\raggedright
\end{minipage} & \begin{minipage}[b]{\linewidth}\raggedright
\end{minipage} & \begin{minipage}[b]{\linewidth}\raggedright
6817000
\end{minipage} & \begin{minipage}[b]{\linewidth}\raggedright
\end{minipage} & \begin{minipage}[b]{\linewidth}\raggedright
\end{minipage} & \begin{minipage}[b]{\linewidth}\raggedright
\end{minipage} \\
\multirow{2}{*}{\begin{minipage}[b]{\linewidth}\raggedright
6000000
\end{minipage}} &
\multirow{2}{*}{\begin{minipage}[b]{\linewidth}\raggedright
\end{minipage}} &
\multirow{2}{*}{\begin{minipage}[b]{\linewidth}\raggedright
\end{minipage}} & &
\multirow{2}{*}{\begin{minipage}[b]{\linewidth}\raggedright
\end{minipage}} &
\multirow{2}{*}{\begin{minipage}[b]{\linewidth}\raggedright
5512000
\end{minipage}} &
\multirow{2}{*}{\begin{minipage}[b]{\linewidth}\raggedright
\end{minipage}} &
\multirow{2}{*}{\begin{minipage}[b]{\linewidth}\raggedright
6203300
\end{minipage}} &
\multirow{2}{*}{\begin{minipage}[b]{\linewidth}\raggedright
\end{minipage}} &
\multirow{2}{*}{\begin{minipage}[b]{\linewidth}\raggedright
Total
\end{minipage}} \\
& & & \begin{minipage}[b]{\linewidth}\raggedright
\end{minipage} \\
\midrule\noalign{}

\bottomrule\noalign{}

4000000 & \begin{minipage}[t]{\linewidth}\raggedright
2782000
\end{minipage} & 3718000 & & 3770000 & & & &
\begin{minipage}[t]{\linewidth}\raggedright
3744000
\end{minipage} & \\
\end{tabular}

2000000

0

\begin{tabular}{@{}
  >{\raggedright\arraybackslash}p{(\columnwidth - 18\tabcolsep) * \real{0.1000}}
  >{\raggedright\arraybackslash}p{(\columnwidth - 18\tabcolsep) * \real{0.1000}}
  >{\raggedright\arraybackslash}p{(\columnwidth - 18\tabcolsep) * \real{0.1000}}
  >{\raggedright\arraybackslash}p{(\columnwidth - 18\tabcolsep) * \real{0.1000}}
  >{\raggedright\arraybackslash}p{(\columnwidth - 18\tabcolsep) * \real{0.1000}}
  >{\raggedright\arraybackslash}p{(\columnwidth - 18\tabcolsep) * \real{0.1000}}
  >{\raggedright\arraybackslash}p{(\columnwidth - 18\tabcolsep) * \real{0.1000}}
  >{\raggedright\arraybackslash}p{(\columnwidth - 18\tabcolsep) * \real{0.1000}}
  >{\raggedright\arraybackslash}p{(\columnwidth - 18\tabcolsep) * \real{0.1000}}
  >{\raggedright\arraybackslash}p{(\columnwidth - 18\tabcolsep) * \real{0.1000}}@{}}
\toprule\noalign{}
\begin{minipage}[b]{\linewidth}\raggedright
P01
\end{minipage} & \begin{minipage}[b]{\linewidth}\raggedright
P02
\end{minipage} & \begin{minipage}[b]{\linewidth}\raggedright
P03
\end{minipage} & \begin{minipage}[b]{\linewidth}\raggedright
P04
\end{minipage} & \begin{minipage}[b]{\linewidth}\raggedright
P05
\end{minipage} & \begin{minipage}[b]{\linewidth}\raggedright
P06
\end{minipage} & \begin{minipage}[b]{\linewidth}\raggedright
P07
\end{minipage} & \begin{minipage}[b]{\linewidth}\raggedright
P08
\end{minipage} & \begin{minipage}[b]{\linewidth}\raggedright
P09
\end{minipage} & \begin{minipage}[b]{\linewidth}\raggedright
P10
\end{minipage} \\
\midrule\noalign{}

\bottomrule\noalign{}

\end{tabular} \\
\midrule\noalign{}

\bottomrule\noalign{}

\end{tabular}

En croisant cette analyse avec celle des quantités, on peut tirer des
conclusions stratégiques importantes :

\begin{itemize}
\item \textbf{Stratégie Volume vs. Valeur :}

\end{itemize}
o \textbf{P02 et P10} sont vos produits de \textbf{flux} (grande
consommation, faible prix unitaire, rotation rapide).

o \textbf{P04, P07 et P08} sont vos produits de \textbf{spécialité}
(haute valeur unitaire, probablement plus complexes à produire ou
destinés à des besoins spécifiques).

\begin{itemize}
\item \textbf{Risque Client/Produit :} Notre rentabilité est mieux répartie que votre production. Si la demande pour P02 chute,
\end{itemize}
l\textquotesingle impact sur votre CA sera moins violent que
l\textquotesingle impact sur votre charge de travail en usine, car les
autres produits soutiennent efficacement la structure financière.

\begin{itemize}
\item \textbf{Analyse du Prix Unitaire :} On constate un écart massif. Par exemple, le produit \textbf{P01} (2,7 millions de CA pour seulement 107
\end{itemize}
unités) a un prix unitaire environ \textbf{650 fois plus élevé} que le
produit \textbf{P02}.
\end{minipage} \\
\midrule\noalign{}
\endhead
\bottomrule\noalign{}
\endlastfoot
\end{longtable}

\begin{center}
\fbox{\begin{minipage}{0.9\linewidth}\centering
\textbf{[PLACEHOLDER: Add image here]}
 \textit{Original image path: media/image23.png}
\end{minipage}}
\end{center}

\begin{center}
\fbox{\begin{minipage}{0.9\linewidth}\centering
\textbf{[PLACEHOLDER: Add image here]}
 \textit{Original image path: media/image4.png}
\end{minipage}}
\end{center}

\begin{center}
\fbox{\begin{minipage}{0.9\linewidth}\centering
\textbf{[PLACEHOLDER: Add image here]}
 \textit{Original image path: media/image14.png}
\end{minipage}}
\end{center}

\begin{center}
\fbox{\begin{minipage}{0.9\linewidth}\centering
\textbf{[PLACEHOLDER: Add image here]}
 \textit{Original image path: media/image15.png}
\end{minipage}}
\end{center}

\begin{center}
\fbox{\begin{minipage}{0.9\linewidth}\centering
\textbf{[PLACEHOLDER: Add image here]}
 \textit{Original image path: media/image16.png}
\end{minipage}}
\end{center}

\begin{center}
\fbox{\begin{minipage}{0.9\linewidth}\centering
\textbf{[PLACEHOLDER: Add image here]}
 \textit{Original image path: media/image17.png}
\end{minipage}}
\end{center}

\begin{center}
\fbox{\begin{minipage}{0.9\linewidth}\centering
\textbf{[PLACEHOLDER: Add image here]}
 \textit{Original image path: media/image18.png}
\end{minipage}}
\end{center}

\begin{center}
\fbox{\begin{minipage}{0.9\linewidth}\centering
\textbf{[PLACEHOLDER: Add image here]}
 \textit{Original image path: media/image19.png}
\end{minipage}}
\end{center}

\begin{center}
\fbox{\begin{minipage}{0.9\linewidth}\centering
\textbf{[PLACEHOLDER: Add image here]}
 \textit{Original image path: media/image5.png}
\end{minipage}}
\end{center}

\begin{center}
\fbox{\begin{minipage}{0.9\linewidth}\centering
\textbf{[PLACEHOLDER: Add image here]}
 \textit{Original image path: media/image20.png}
\end{minipage}}
\end{center}

\begin{longtable}[]{@{}
  >{\raggedright\arraybackslash}p{(\columnwidth - 0\tabcolsep) * \real{1.0000}}@{}}
\toprule\noalign{}
\begin{minipage}[b]{\linewidth}\raggedright
Pareto

\begin{tabular}{@{}
  >{\raggedright\arraybackslash}p{(\columnwidth - 22\tabcolsep) * \real{0.0833}}
  >{\raggedright\arraybackslash}p{(\columnwidth - 22\tabcolsep) * \real{0.0833}}
  >{\raggedright\arraybackslash}p{(\columnwidth - 22\tabcolsep) * \real{0.0833}}
  >{\raggedright\arraybackslash}p{(\columnwidth - 22\tabcolsep) * \real{0.0833}}
  >{\raggedright\arraybackslash}p{(\columnwidth - 22\tabcolsep) * \real{0.0833}}
  >{\raggedright\arraybackslash}p{(\columnwidth - 22\tabcolsep) * \real{0.0833}}
  >{\raggedright\arraybackslash}p{(\columnwidth - 22\tabcolsep) * \real{0.0833}}
  >{\raggedright\arraybackslash}p{(\columnwidth - 22\tabcolsep) * \real{0.0833}}
  >{\raggedright\arraybackslash}p{(\columnwidth - 22\tabcolsep) * \real{0.0833}}
  >{\raggedright\arraybackslash}p{(\columnwidth - 22\tabcolsep) * \real{0.0833}}
  >{\raggedright\arraybackslash}p{(\columnwidth - 22\tabcolsep) * \real{0.0833}}
  >{\raggedright\arraybackslash}p{(\columnwidth - 22\tabcolsep) * \real{0.0833}}@{}}
\toprule\noalign{}
\begin{minipage}[b]{\linewidth}\raggedright
12000000
\end{minipage} &
\multirow{11}{*}{\begin{minipage}[b]{\linewidth}\raggedright
P02
\end{minipage}} &
\multirow{11}{*}{\begin{minipage}[b]{\linewidth}\raggedright
P10
\end{minipage}} &
\multirow{11}{*}{\begin{minipage}[b]{\linewidth}\raggedright
P04
\end{minipage}} &
\multirow{11}{*}{\begin{minipage}[b]{\linewidth}\raggedright
P07
\end{minipage}} &
\multirow{11}{*}{\begin{minipage}[b]{\linewidth}\raggedright
P08
\end{minipage}} &
\multirow{11}{*}{\begin{minipage}[b]{\linewidth}\raggedright
P06
\end{minipage}} &
\multirow{11}{*}{\begin{minipage}[b]{\linewidth}\raggedright
P05
\end{minipage}} &
\multirow{11}{*}{\begin{minipage}[b]{\linewidth}\raggedright
P09
\end{minipage}} &
\multirow{11}{*}{\begin{minipage}[b]{\linewidth}\raggedright
P03
\end{minipage}} &
\multirow{11}{*}{\begin{minipage}[b]{\linewidth}\raggedright
P01
\end{minipage}} & \begin{minipage}[b]{\linewidth}\raggedright
100\%
\end{minipage} \\
\multirow{2}{*}{\begin{minipage}[b]{\linewidth}\raggedright
10000000
\end{minipage}} & & & & & & & & & & &
\begin{minipage}[b]{\linewidth}\raggedright
90\%
\end{minipage} \\
& & & & & & & & & & & \begin{minipage}[b]{\linewidth}\raggedright
80\%
\end{minipage} \\
\begin{minipage}[b]{\linewidth}\raggedright
8000000
\end{minipage} & & & & & & & & & & &
\begin{minipage}[b]{\linewidth}\raggedright
70\%
\end{minipage} \\
\multirow{2}{*}{\begin{minipage}[b]{\linewidth}\raggedright
6000000
\end{minipage}} & & & & & & & & & & &
\begin{minipage}[b]{\linewidth}\raggedright
60\%
\end{minipage} \\
& & & & & & & & & & & \begin{minipage}[b]{\linewidth}\raggedright
50\%
\end{minipage} \\
\multirow{2}{*}{\begin{minipage}[b]{\linewidth}\raggedright
4000000
\end{minipage}} & & & & & & & & & & &
\begin{minipage}[b]{\linewidth}\raggedright
40\%
\end{minipage} \\
& & & & & & & & & & & \begin{minipage}[b]{\linewidth}\raggedright
30\%
\end{minipage} \\
\begin{minipage}[b]{\linewidth}\raggedright
2000000
\end{minipage} & & & & & & & & & & &
\begin{minipage}[b]{\linewidth}\raggedright
20\%
\end{minipage} \\
\multirow{2}{*}{\begin{minipage}[b]{\linewidth}\raggedright
0
\end{minipage}} & & & & & & & & & & &
\begin{minipage}[b]{\linewidth}\raggedright
10\%
\end{minipage} \\
& & & & & & & & & & & \begin{minipage}[b]{\linewidth}\raggedright
0\%
\end{minipage} \\
\midrule\noalign{}

\bottomrule\noalign{}

\end{tabular}

Priorité de suivi : Concentrez les efforts
d\textquotesingle optimisation des coûts et de qualité sur les produits
de la Zone A (P02, P10) et le début de la Zone B (P04, P07). Une
réduction de coût de 5\% sur P02 a beaucoup plus
d\textquotesingle impact sur la rentabilité globale
qu\textquotesingle une réduction de 20\% sur P01.
\end{minipage} \\
\midrule\noalign{}
\endhead
\bottomrule\noalign{}
\endlastfoot
\end{longtable}

\begin{longtable}[]{@{}
  >{\raggedright\arraybackslash}p{(\columnwidth - 0\tabcolsep) * \real{1.0000}}@{}}
\toprule\noalign{}
\begin{minipage}[b]{\linewidth}\raggedright
\begin{center}
\fbox{\begin{minipage}{0.9\linewidth}\centering
\textbf{[PLACEHOLDER: Add image here]}
 \textit{Original image path: media/image4.png}
\end{minipage}}
\end{center}

\begin{center}
\fbox{\begin{minipage}{0.9\linewidth}\centering
\textbf{[PLACEHOLDER: Add image here]}
 \textit{Original image path: media/image5.png}
\end{minipage}}
\end{center}

Après avoir réalisé des graphiques quantitatifs on peut passer a la
partie d'implantation par cela on trouve que les produit ayant le plus
grand impact sont P02 et P10 c'est pour cela on a décidé de zoomer et
voir la partie temporelle :

\begin{center}
\fbox{\begin{minipage}{0.9\linewidth}\centering
\textbf{[PLACEHOLDER: Add image here]}
 \textit{Original image path: media/image24.png}
\end{minipage}}
\end{center}

\begin{center}
\fbox{\begin{minipage}{0.9\linewidth}\centering
\textbf{[PLACEHOLDER: Add image here]}
 \textit{Original image path: media/image25.png}
\end{minipage}}
\end{center}

\end{minipage} \\
\midrule\noalign{}
\endhead
\bottomrule\noalign{}
\endlastfoot
\end{longtable}

\begin{longtable}[]{@{}
  >{\raggedright\arraybackslash}p{(\columnwidth - 0\tabcolsep) * \real{1.0000}}@{}}
\toprule\noalign{}
\begin{minipage}[b]{\linewidth}\raggedright
\begin{center}
\fbox{\begin{minipage}{0.9\linewidth}\centering
\textbf{[PLACEHOLDER: Add image here]}
 \textit{Original image path: media/image4.png}
\end{minipage}}
\end{center}

\begin{center}
\fbox{\begin{minipage}{0.9\linewidth}\centering
\textbf{[PLACEHOLDER: Add image here]}
 \textit{Original image path: media/image5.png}
\end{minipage}}
\end{center}

Pour essayer de resume

MLT : manufacturing Lead time\\
Temps de passage\\
\hl{Temps à valeur ajouter}\\
\hl{Temps a non-valeur ajouter}\\
MLT = \hl{Temps de machine + Temps de chargement de format + Temps des
en cours} + \hl{Temps Attente} \hl{+ Temps déplacement}\\
TVA/MLT = TVA/(TVA+TNVA) SI = 1 LEAN SI \textgreater1 NON-APPLICATION DU
LEAN CAR TNVA TROP IMPORTANT
\end{minipage} \\
\midrule\noalign{}
\endhead
\bottomrule\noalign{}
\endlastfoot
\end{longtable}

\end{document}
